\begin{titlepage}
  \begin{center}

  {\Huge DC\_BLOCK\_RAM}

  \vspace{25mm}

  \includegraphics[width=0.90\textwidth,height=\textheight,keepaspectratio]{img/AFRL.png}

  \vspace{25mm}

  \today

  \vspace{15mm}

  {\Large Jay Convertino}

  \end{center}
\end{titlepage}

\tableofcontents

\newpage

\section{Usage}

\subsection{Introduction}

\par
Dual clock block RAM for any FPGA target. Includes a byte enable for selecting bytes to write from the bus.

\subsection{Dependencies}

\par
The following are the dependencies of the cores.

\begin{itemize}
  \item fusesoc 2.X
  \item iverilog (simulation)
  \item cocotb (simulation)
\end{itemize}

\subsubsection{fusesoc\_info Depenecies}
\begin{itemize}
\item dep
	\begin{itemize}
	\item AFRL:utility:helper:1.0.0
	\end{itemize}
\item dep\_tb
	\begin{itemize}
	\item AFRL:simulation:clock\_stimulator
	\item AFRL:utility:sim\_helper
	\end{itemize}
\end{itemize}


\subsection{In a Project}
\par
Connect the device using the read write signals see \ref{Module Documentation} for details

\section{Architecture}
\par
This core is made up of a single module.
\begin{itemize}
  \item \textbf{ft245\_sync\_to\_axis} Interface AXIS to F245 device (see core for documentation).
\end{itemize}

\par
This core has 2 always blocks that are sensitive to the positive clock edge.

\begin{itemize}
\item \textbf{Produce Data} Takes write input data and stores it in RAM at a specified address. BE will filter out bytes if the corresponding bits not set to active high.
\item \textbf{Consume Data} Read data from RAM at a specified address and output over read interface.
\end{itemize}

Please see \ref{Module Documentation} for information on read/write interface ports.

\section{Building}

\par
The DC block RAM is written in Verilog 2001. It should synthesize in any modern FPGA software. The core comes as a fusesoc packaged core and can be included in any other core. Be sure to make sure you have meet the dependencies listed in the previous section. Linting is performed by verible using the lint target.

\subsection{fusesoc}
\par
Fusesoc is a system for building FPGA software without relying on the internal project management of the tool. Avoiding vendor lock in to Vivado or Quartus.
These cores, when included in a project, can be easily integrated and targets created based upon the end developer needs. The core by itself is not a part of
a system and should be integrated into a fusesoc based system. Simulations are setup to use fusesoc and are a part of its targets.

\subsection{Source Files}

\subsubsection{fusesoc\_info File List}
\begin{itemize}
\item src
	\begin{itemize}
	\item src/dc\_block\_ram.v
	\end{itemize}
\item tb
	\begin{itemize}
	\item {'tb/tb\_fifo.v': {'file\_type': 'verilogSource'}}
	\end{itemize}
\end{itemize}


\subsection{Targets}

\subsubsection{fusesoc\_info Targets}
\begin{itemize}
\item default
	\begin{itemize}
	\item[$\space$] Info: Default for IP intergration.
	\end{itemize}
\item lint
	\begin{itemize}
	\item[$\space$] Info: Lint with Verible
	\end{itemize}
\item sim
	\begin{itemize}
	\item[$\space$] Info: Default for IP intergration.
	\end{itemize}
\end{itemize}


\subsection{Directory Guide}

\par
Below highlights important folders from the root of the directory.

\begin{enumerate}
  \item \textbf{docs} Contains all documentation related to this project.
    \begin{itemize}
      \item \textbf{manual} Contains user manual and github page that are generated from the latex sources.
    \end{itemize}
  \item \textbf{src} Contains source files for the core
  \item \textbf{tb} Contains test bench files for iverilog and cocotb
\end{enumerate}

\newpage

\section{Simulation}
\par
There are a few different simulations that can be run for this core.

\subsection{iverilog}
\par
iverilog is used for simple test benches for quick verification, visually, of the core.

\subsection{cocotb}
\par
This feature is not implemented for this core.

\newpage

\section{Module Documentation} \label{Module Documentation}

\begin{itemize}
\item \textbf{dc\_block\_ram} Generic dual clock block RAM\\
\end{itemize}
The next sections document the module.

